% Folder 93, batch 03 — archivists' pages 41–60.
% Pass: Opus 5.5 (claude-opus-5-5), 2026-10-10 — first pass, unchecked against the pages by a human.
%
% Contents: continuation of his annotated typescript on Dieudonné theory
% (pp. 41–47: filtered F-V-crystals, deformation theorems for BT groups and
% abelian schemes, admissible F-V-crystals, comparison of the two theories),
% a further typescript of "rabiots" and problems (pp. 48, 50, 52–53: isogroups
% of BT, Pisa notes VII 10.8.2, Tate's theorem, Problems 4–6), and manuscript
% notes "chi et chi_!" on Euler characteristics with and without compact
% supports (pp. 55, 57, 59). Skipped without trace: pp. 49 and 51 (Bourbaki
% draft n° 393 on Banach algebras, versos), p. 54 (cover, his words used as
% heading), pp. 56, 58, 60 (administrative typescript versos).
\documentclass[11pt,a4paper]{article}
\input{../preamble/grothendieck}

\folder{93}
\batch{3}
\pages{41}{60}
\foldertitle{Connexions de Gauss-Manin et équations de Picard-Fuchs. Opération de Cartier : copies de tapuscrit annoté (s.d.), tiré à part (1981), notes manuscrites (s.d.).}
\dating{[1972]-1981}
\watermark{Édition de démonstration}

\begin{document}

\page{41}
\note{la page est une copie d'un tapuscrit, annotée à la main ; elle prend un exposé en cours (la filtration (*) et le problème 1 renvoient à ce qui précède la tranche)}

Appelons F-V-\struck{cristal} \emph{\uncertain{cristal} filtré} \add{sur $S$} (par abus de langage) un F-V-cristal en Mod.\ loc.\ lib.\ sur $S$ muni \struck{de filtration} de la donnée d'une filtration de $M_S$ par une sous-module localement facteur direct. On trouve donc en fait un foncteur
\[
  \mathrm{BT}(S) \longrightarrow \text{F-V-cris fil}\,(S) ,
\]
compatible aux images inverses.

\emph{Remarque.} \struck{La} filtration (*) doit être considérée comme étant l'équivalent d'une filtration «~de Hodge~» sur une cohomologie de De~Rham \add{relative} en dimension~1. \struck{Pour expliciter ce point, partons par exemple d'un schéma abélien $A$ sur $S$, et prenons $G = A(\infty) = \varinjlim A(n)$. On peut alors} \add{sera explicité plus bas.}

Soit maintenant $S'$ un épaississement à p.div.\ de $S$. On se propose de trouver tous les prolongements possibles (à isom.\ près) du groupe de BT $G$ donné sur $S$ en un groupe de BT $G'$ sur $S'$. On \struck{trouve} prouvera que si à tout tel $G'$, on associe la filtration correspondante de $\mathbb{D}^*(G')_{S'} = \mathbb{D}^*(G)_{S'} = \mathbb{D}^*(G_0)_{S'}$ (cf.\ \note{renvoi laissé en blanc} pour le fait que $S'$ peut être envisagé comme un objet du site cristallin absolu de $S_0$), qui prolonge la filtration qu'on a sur $\mathbb{D}^*(G)_S = \mathbb{D}^*(G_0)_S = \mathbb{D}^*(G')|S$, on trouve une bijection entre l'ensemble des classes d'isomorphie cherchés et l'ensemble des filtrations qui prolongent. De façon plus précise, on a\,: \struck{une équivalence de catégories}

\emph{Théorème de déformation pour les groupes de BT}\,: Soient $S$ un schéma où $p$ est loc.\ nilp., $S'$ un épaississement à p.div.\ de $S$, Considérons le foncteur canonique
\[
  \mathrm{BT}(S') \Longrightarrow
\]
catégorie des \struck{triples} \add{paires} d'un groupe de BT $G$ sur $S$, et d'une sous-module localement facteur direct de $\mathbb{D}^*(G)_{S'}$ qui prolonge le sous-module $\mathrm{Fil}^1\mathbb{D}^*(G)_S = \underline{\omega}_G$ de $\mathbb{D}^*(G)_S$.
\note{la machine n'avait pas le $\mathbb{D}$ ; le tapuscrit tape D et le souligne à la main, ici et dans toute la tranche}

Ce foncteur est une \emph{équivalence de catégories}. \add{si $S'$ est à p.\ div.\ nil\uncertain{p.}, \ill{}}

\marginal{Remarques\,: si on se borne à des $G$, $G'$ qui sont infinitésimaux, on \uncertain{a} \ill{}}
\note{l'ajout manuscrit court sur deux lignes, au-dessus et au-dessous d'un trait ; la fin, coupée par le bord, ne se lit pas}

Si le \struck{question} problème~1 avait une réponse affirmative, on en conclurait une description de la catégorie des groupes de BT sur $S$ en termes de cristaux de Dieudonné filtrés sur $S$, en appliquant le théorème de

\page{42}
déformation au\struck{x} cas du couple $(S_0, S)$, \struck{\ill{}} notant que l'idéal $p\underline{O}_S$ est bien un idéal à puissances divisées. On notera d'ailleurs que le théorème de déformations dans le cas \struck{envisagé} général $(S,S')$ envisagé dans son énoncé est en fait équivalent au cas particulier d'un couple $(S_0,S)$, comme il est formel via $\mathrm{Crisloclib}(S) \xrightarrow{\approx} \mathrm{Crisloclib}(S_0)$.

\add{\emph{Problème~3}\,: Trouver variante du th.\ de déformation pr les p-Groupes finis loc.\ libres (j'ai un candidat).}

c) \emph{F-V-cristaux admissibles en car.\ p.} Revenons au cas d'un schéma de car.\ $p$. Le cristal de Dieudonné $\mathbb{D}^*(G_0)$ étant muni d'une filtration canonique (de $\mathbb{D}^*(G_0)_{S_0}$), on pourrait craindre que le problème~1 est mal posé, et qu'il faudrait envisager le foncteur comme étant à valeurs dans des cristaux de Dieudonné filtrés. En fait, dans le cas envisagé ici, il se trouve que la filtration est uniquement déterminée en termes de la structure de cristal de Dieudonné, comme on va l'expliquer maintenant.

Considérons le morphisme de topos annelés (qui aura été défini déjà précédemment, dans le par.\ d'exemples de cristaux, pour les relations entre cristaux et vecteurs de Witt)
\[
  \phi : S_{0\,\mathrm{cris}/\underline{F}_p} \overset{\varepsilon}{\longrightarrow} S_0 .
\]
Soit $M = \mathbb{D}^*(G_0)$, et soit $M_0$ sa restriction au site cristallin rel.\ à $\underline{F}_p$. Il se trouve que \struck{la suite} d'après $FV = p$, $VF = p$, les composés dans la suite
\[
  M_0^{(p)} \xrightarrow{F_{M_0}} M_0 \xrightarrow{V_{M_0}} M_0^{(p)} \xrightarrow{F_{M_0}} M_0
\]
sont nuls\,; en fait, on prouvera mieux, à savoir que \add{pour tout cristal de Dieud.\ $M$,} la suite précédente est exacte (terme à terme sur tout épaississement à p.div.\ de car.\,p de $S_0$) et on obtient ainsi des sous-cristaux en modules loc.\ libres
\[
  \operatorname{Ker} V_{M_0} = \operatorname{Im} F_{M_0} \subset M_0 , \qquad
  \operatorname{Ker} F_{M_0} = \operatorname{Im} V_{M_0} \subset M_0^{(p)} .
\]
Ceci dit, \struck{le sous-Module $\mathrm{Fil}^1$} notons que $M_{S_0} = \varepsilon^*(M_0)$, donc
\[
  \phi^*(M_{S_0}) \simeq M_0^{(p)}
\]
donc on a
\[
  \phi^*(\mathrm{Fil}^1(M_{S_0})) \subset M_0^{(p)} .
\]
Ceci posé, on montre qu'on a

\page{43}
\[
  \phi^*(\mathrm{Fil}^1(M_{S_0})) = \operatorname{Ker} F_{M_0} \;(= \operatorname{Im} V_{M_0}) \subset M_0^{(p)} .
\]
D'autre part, on montre que $\phi$ a des propriétés de fidélité telles que la relation précédente détermine sans ambiguité le sous-Module localement facteur direct $\mathrm{Fil}^1$ de $M_{S_0}$.

On dira qu'un \add{F-V-}cristal en modules loc.\ lib.\ $M_0$ sur le site cristallin de $S_0$ au dessus de $\underline{F}_p$ est \emph{admissible} s'il existe un sous-module loc.\ facteur direct $\mathrm{Fil}^1$ de $M_{0S}$ tel que l'on ait la relation précédente\,; \struck{Ces} \add{Ce} $\mathrm{Fil}^1$ est unique et \struck{déterminé} détermine une filtration qu'on appellera la \emph{filt.\ canonique}. On dit de même qu'un cristal de Dieudonné \add{$M$} sur $S_0$ est admissible si sa restriction $M_0$ \struck{sur} au site cristallin rel.\ à $\underline{F}_p$ est admissible.
\note{la phrase « détermine une filtration … canonique. On » est tapée en tête de la page, entourée, et un trait la renvoie après « unique et »}
On voit donc que si $M$ est le cristal de Dieudonné associé à un groupe de BT sur $S_0$, $M$ est admissible, et la filtration de $M_{S_0}$ envisagée dans b) est la filtration canonique définie précédemment.

\marginal{d) Passage à la limite\,: \uncertain{formule} anneau $A$ \uncertain{séparé et} complet pour top.\ \emph{p-adique}. e) Th\ill{} \ill{} et th.\ de Tate. f) Dualité.}

2)\note{le chiffre est surchargé ; « 2) » s'accorde avec le « 3) » de la p.~47} \emph{Théorie de déformations pour les schémas abéliens.}

\add{a)} Dans cette théorie, il n'y a plus de nombre premier privilégié, par contre il faut travailler avec le site cristallin \emph{nilpotent}. On trouve un foncteur \add{contravariant} de la catégorie des schémas abéliens sur le schéma $S$ vers les \struck{faisceaux} cristaux sur le site cristallin nilpotent de $S$
\[
  \mathrm{Schémab}(S) \longrightarrow \mathrm{Crisloclib}_{\mathrm{nilp}}(S)
\]
qui peut être défini par exemple comme
\[
  (*) \qquad \mathbb{D}^* : A \longmapsto R^1 f_{A\mathrm{cris}*}(\underline{O}_{A_{\mathrm{cris}}})
\]
où on considère
\[
  f_{\mathrm{cris}} : A_{\mathrm{crisnilp}} \longrightarrow S_{\mathrm{crisnilp}}
\]
induit par $S$. Par une variante du théorème de changement de base, on voit que
\[
  \mathbb{D}^*(A)_S = \text{\struck{\ill{}}}\; \underline{H}^1_{\mathrm{DR}}(A/S) \overset{\mathrm{dfn}}{=} R^1 f_{A*}(\Omega^{\bullet}_{A/S}) .
\]
Pour calculer $\mathbb{D}^*(A)_{S'}$, $S'$ épaississement à p.div.\ loc.\ nilpotentes de $S$, on aura pour tout prolongement \struck{canonique} de $A$ en un schéma abélien $A'$ sur $S'$

\page{44}
\[
  \mathbb{D}^*(A)_{S'} = \text{\struck{\ill{}}}\; \underline{H}^1_{\mathrm{DR}}(A'/S') .
\]
On notera que d'après un théorème connu (cf.\ livre de Mumford Geometric Invariant theory) il existe toujours un $A'$ pour $S$ affine\,; le fait que \struck{\ill{}} \add{$H^1(A'/S')$} ne dépende, à isomorphisme près, pas du prolongement choisi de $A$, peut être considéré ici comme conséquence du tapis de la cohomologie cristalline.

En fait, utilisant la filtration de \struck{De Rham} \add{Hodge} sur la cohomologie de De~Rham, on trouve une filtration sur $\mathbb{D}^*(A)_S$, savoir
\[
  0 \longrightarrow R^0 f_*(\underline{\Omega}^1_{A/S}) \longrightarrow \underline{H}^1_{\mathrm{DR}}(A/S) \longrightarrow R^1 f_*(\underline{O}_S) \longrightarrow 0 ,
\]
qu'on peut interpréter comme
\[
  (**) \qquad 0 \longrightarrow \underset{\substack{\parallel \\ \underline{\omega}_A}}{\underline{\check{t}}_A} \longrightarrow \mathbb{D}^*(A)_S \longrightarrow \underline{t}_{A^*} \longrightarrow 0 .
\]
\note{l'égalité $\check{t}_A = \underline{\omega}_A$ est écrite à la main sous la flèche}
Le foncteur (*) et la suite exacte (**) sont fonctoriels en $A$ et compatibles avec image inverse. On peut \struck{envisager} considérer que le foncteur (*) se factorise en
\[
  \mathbb{D}^* : \mathrm{Schémab}(S) \longrightarrow \mathrm{Crisloclib}_{\mathrm{nilp}}\,\mathrm{fil}(S)
\]
vers les cristaux en modules localement libres «~filtrés~» (en deux crans localement libres). Utilisant cette factorisation, on trouve que pour un objet $S'$ du site cristallin nilpotent de $S$, la donnée d'un prolongement de $A$ en un schéma \add{abélien} $A'$ sur $S'$ revient à la donnée d'un prolongement de la filtration canonique (\struck{(**)}) de $\mathbb{D}^*(A)_S$ en une filtration (à quotients localement libres) de $\mathbb{D}^*(A)_{S'}$. Plus précisément, on a le

\emph{Théorème de déformation pour les schémas abéliens.} Soient $S$ un schéma, $S'$ un voisinage à puissances div.\ loc.\ nilp.\ de $S$, alors le foncteur naturel
\[
  \mathrm{Schémab}(S') \longrightarrow
\]
catégorie des couples $(A, \mathrm{Fil}^1)$ d'un schéma abélien $A$ sur $S$ et d'un sous-module localement facteur direct $\mathrm{Fil}^1$ \add{de $\mathbb{D}^*(A)_{S'}$} prolongeant $\mathrm{Fil}^1\mathbb{D}^*(A)_S \simeq \underline{\omega}_A$

est une équivalence de catégories.

\add{b)} On peut donner une deuxième interprétation du foncteur $\mathbb{D}^*$ sur les

\page{45}
schémas abéliens, \add{permettant de définir un foncteur qu.\,inverse du foncteur du théorème.} Pour ceci, on considère pour tout schéma abélien $A$ sur $S$ son extension vectorielle universelle, qui est une extension
\[
  0 \longrightarrow \underline{\check{t}}_{A^*} \longrightarrow E(A) \longrightarrow A \longrightarrow 0 .
\]
\note{la phrase « permettant … du théorème. » est tapée au-dessus de la ligne et entourée, avec un renvoi après « schémas abéliens, »}
Il se trouve qu'on peut trouver un cristal en groupes lisses canonique sur $S$, soit $\underline{E}(A)$, tel que l'on ait
\[
  E(A) \simeq \underline{E}(A)_S .
\]
La définition est telle que \struck{pour tout le foncteur} $A \mapsto \underline{E}(A)$ définisse un foncteur
\[
  \mathrm{Schémab}(S) \longrightarrow \mathrm{Cris\ Groupes\ lisses}_{\mathrm{nilp}}(S)
\]
compatible aux changements de base, l'isomorphisme précédent étant également fonctoriel et compatible aux changements de base. Il s'ensuit donc que l'on peut définir $\underline{E}(A)$ comme $S' \mapsto \underline{E}(A)_{S'}$, avec
\[
  \underline{E}(A)_{S'} = E(A') ,
\]
$A'$ désignant un schéma abélien qui prolonge $A$ sur $S'$, \struck{\ill{}} à charge de définir \struck{des} \add{un système transitif d'}isomorphismes canoniques entre les groupes obtenus pour des prolongements différents. Dans le cas où $S$ est de caractéristique nulle (donc plus question de puissances divisées\,!) cela est particulièrement simple\,; on prouve qu'il existe à isomorphisme \emph{unique} près un seul schéma en groupes lisses $E'$ sur l'épaississement $S'$ de $S$ qui prolonge $E$ (NB il suffirait même d'un nilépaississement, comme on voit par passage à la limite à partir du cas noethérien)\,: c'est le résultat déjà signalé dans ma vieille lettre à Tate. Dans le cas général, c'est plus délicat, et peut se faire par «~la méthode de l'exponentielle~», qui sera exposée dans le contexte analogue des groupes de BT dans la suite du séminaire. On peut aussi définir $\underline{E}(A^*)$ par
\[
  \underline{E}(A^*) = R^1 f_{A\mathrm{cris}*}(\underline{G}_{m\,A_{\mathrm{cris}}}) ,
\]
où $A_{\mathrm{cris}}$ désigne le topos cristallin \struck{abs} nilpotent absolu (relativement à $\underline{Z}$)\,: cette méthode peut également s'adapter au cas des groupes de BT.

Quoi qu'il en soit de la méthode utilisée pour définir $\underline{E}(A)$, on

\page{46}
peut énoncer qu'il y a un isomorphisme canonique (fonctoriel, compatible aux extensions de la base)
\[
  \text{\struck{\ill{}}}\; \mathbb{D}^*(A) = \underline{\mathrm{Lie}}(\underline{E}(A^*)) ,
\]
la structure d'extension \quad de $\mathbb{D}^*(A)$ n'étant autre que la structure \add{sur Lie} déduite de la structure d'extension \quad de $\underline{E}(A)_S = E(A)$.
\note{deux blancs sont laissés dans la ligne du tapuscrit, après chacun des deux « structure d'extension »}

Ceci posé, revenant aux conditions du théorème de déformations, on voit, par un sorite sur l'exponentielle sur lequel nous reviendrons, que la donnée d'un prolongement $\mathrm{Fil}^1 \subset \mathbb{D}^*(A)_{S'} = \underline{\mathrm{Lie}}(\underline{E}(A^*)_{S'})$ localement facteur direct prolongeant $\underline{\mathrm{Lie}}(\underline{\check{t}}_A) = \underline{\check{t}}_A \subset \mathbb{D}^*(A)_S = \underline{\mathrm{Lie}}(\underline{E}(A^*)_{S'}|S)$ équivaut à la donnée d'un sous-groupe \add{lisse} à structure vectorielle de $\underline{E}(A^*)_{S'}$ qui prolonge le sous-groupe $\underline{\check{t}}_A$ de $\underline{E}(A^*)_S$. Si $L$ est ce sous-groupe, alors il est immédiat que \struck{\ill{}}
\[
  \underline{E}(A^*)_{S'}/L
\]
est un schéma abélien sur $S'$, dont le schéma abélien dual est le $A'$ cherché. (NB Il serait plus naturel ici de travailler avec \struck{$\mathbb{D}^*(A) = \mathbb{D}(A)$}
\[
  \mathbb{D}_*(A) = \underline{\mathrm{Lie}}\,\underline{E}(A) \xrightarrow[\sim]{\text{isom can}} \mathbb{D}^*(A^*)
\]
plutôt qu'avec $\mathbb{D}^*(A)$ pour énoncer le théorème de déformations\,; l'équivalence des deux points de vue provient de l'accouplement parfait
\[
  \mathbb{D}^*(A) \otimes \mathbb{D}^*(A^*) \longrightarrow \underline{O}_{S_{\mathrm{crisnilp}}}
\]
compatible avec les filtrations.)

\emph{Remarque.} En théorie des groupes de BT, nous donnerons deux constructions de $\mathbb{D}^*(G)$, inspirées des deux constructions indiquées dans le cas des schémas abéliens\,: l'une «~cohomologique~», l'autre via une «~extension vectorielle universelle~» $\underline{E}(G)$. La construction de $\underline{E}(G)$ peut aussi s'obtenir par deux méthodes, l'une cohomologique, l'autre «~par l'exponentielle~» étant directement adaptée à la démonstration du théorème de déformations.

\page{47}
3) \struck{2) Relations} Relations entre les deux théories (pour schémas abéliens et pour groupes de BT). Supposons que $p$ soit localement nilpotent sur $S$. Travaillons \struck{On suppose} à nouveau avec le site cristallin de Berthelot (pas le nilpotent). Le résultat essentiel sera, pour un schéma abélien variable $A$ sur $S$, un isomorphisme fonctoriel compatible aux changements de base
\[
  \mathbb{D}^*(A(\infty)) \simeq R^1 f_{\mathrm{cris}*}(\underline{O}_{A_{\mathrm{cris}}}) ,
\]
cet isomorphisme induisant, pour les valeurs des deux membres sur $S$, un isomorphisme compatible aux filtrations \quad et \quad (compte tenu des isomorphismes $\underline{t}_{A(\infty)} \simeq \underline{t}_A$ et idoine pour $A^*$). \add{-- C'est en fait en postulant un tel isomorphisme qu'on arrive par voie heuristique à une définition de $\mathbb{D}^*(G)$ pour un groupe de BT $G$ quelconque.}
\note{deux blancs dans la ligne du tapuscrit, de part et d'autre de « et »}

Ceci impliquera en particulier le fait que, pour un \struck{voisinage} épaississement à p.\,div.\ loc.\ nilp.\ \struck{de} $S'$ de $S$, il y a identité de la théorie des prolongements infinitésimaux de $A$, et de $A(\infty)$. De ceci, par dévissage sur l'ordre de nilpotence, on peut déduire par exemple immédiatement le théorème de Serre-Tate déjà annoncé plus haut\,: \add{méthode suivie par Messing ds sa thèse.}

On peut aussi définir une extension vectorielle universelle $E(G)$ d'un groupe de BT, et on trouvera une isomorphisme canonique etc
\[
  E(A(\infty)) \simeq E(A)(\infty)
\]
donnant naissance, par le procédé habituel, à un isomorphisme plus général de cristaux en groupes
\[
  \underline{E}(A(\infty)) \simeq \underline{E}(A)(\infty) \quad \ldots
\]
\note{ce dernier alinéa est encadré au crayon par deux crochets}

\struck{4) Passage à la limite}

\struck{4) Problème\,: trouver le foncteur des F-cristaux filtrés vers modules faisceaux p-adiques.}\note{le mot « modules » est en outre biffé à la machine}

\page{48}
\note{la page est d'un autre tapuscrit (papier jauni) ; elle ouvre une liste de compléments}

Rabiots pour l'exposé précédent\,:

a) Description de la filtration canonique d'un F-V-cristal dans le cas d'une base \add{un anneau \emph{parfait}} de car.\ $p$\,: elle est définie par un sous-module $\mathrm{Fil}^1 M_0$ \struck{($M$ considéré comme module localement libre de type fini sur $W = \underline{W}(A)$)} tel que l'on ait
\[
  \mathrm{Fil}^{1\,(p)} = \operatorname{Ker} F_M = \operatorname{Im} V_M ,
\]
ce qui le définit de façon unique grâce à l'hypothèse sur $A$ (NB tout cristal de Dieudonné est alors «~admissible~»).

b) Isogroupes \add{de BT} (devrait faire un sous-paragraphe). Définition sur une base quelconque\,: on localise \struck{pour $Z$} en remplaçant $\mathrm{Hom}(G,H)$ par $\mathrm{Hom}(G,H) \otimes_{\underline{Z}_p} \underline{Q}_p = \mathrm{Hom} \otimes_{\underline{Z}} \underline{Q}$. (Définition provisoire\,; elle \struck{Proposition} a l'inconvénient de ne pas fournir un champ, et devrait peut-être être reconsidérée~…). On définit de même les isocristaux de Dieudonné, et on obtient donc un foncteur
\[
  \mathrm{IsoBT}(S) \qquad \mathrm{Iso\,F\text{-}V\text{-}cris}(S)
\]
($p$ localement nilpotent sur $S$). On trouve
\note{le mot « Isogroupes » est tapé par-dessus un autre mot ; la flèche manque entre les deux catégories}

\emph{Proposition} (prouvée plus tard)\,: \struck{Supposons $S$ quasi-compact, et soit $S_0$ le sous-schéma défini par l'idéal} Soit $S \to S'$ une immersion nilpotente d'Idéal $J$ annulé par une puissance de $J$, alors le foncteur restriction
\[
  \mathrm{IsoBT}(S') \longrightarrow \mathrm{IsoBT}(S)
\]
est une équivalence de catégories.

\marginal{Th.}
On en conclut l'énoncé VII~10.8.2 de mes notes de Pise. \add{Expliciter le cas d'un anneau de val.\ discrète d'inégales car.} On peut se poser \struck{la question de l'image essentielle} le

\emph{Problème}~4. Quel est l'image essentielle dudit foncteur\,? Suppposant $A$ noethérien séparé et complet pour la topologique p-adique pour fixer les idées, y a-t-il un critère valuatif d'appartenance à l'image essentielle\,? Si $A$ est un anneau de val.\ discrète d'inégales caractéristiques, l'image essentielle est elle un ouvert de Zariski de la grassma-

\page{50}
nienne sur le corps des fractions \struck{de} $L$ de $K$\,?

On notera que même dans le cas où $A = W(k)$, $k$ alg.\ clos, le foncteur en question n'est pas \struck{essentiellem} essentiellement surjectif. Notons le

\emph{Problème}~5. (A anneau de val.\ discrète complet d'inégales car., $k$ parf) Expliciter le foncteur naturel qui va de la sous-catégorie pleine précédente de la catégorie des couples $(M, \mathrm{Fil})$ d'un module de Dieudonné $M$ sur \struck{$W$} $k$ (i.e.\ \struck{sur} module sur $W$) et d'une filtration à deux crans de $M \otimes_W L$, vers les modules galoisiens sur $L$ (représentations vectorielles des dimension finie sur $\underline{Q}_p$ de $\mathrm{Gal}(\overline{L}/L)$) \struck{dans des espaces modules libres de type fini}. Peut-on même définir ce foncteur \struck{sur la sous-catégorie} de façon naturelle sans se borner à cette sous-catégorie pleine, en imposant seulement des conditions de rangs sur le gradué, mettons\,?

\emph{Problème}~6. Revenant à un $A$ comme dans Pb~4, peut-on trouver un foncteur \struck{dans} de la catégorie des couples $(M, \mathrm{Fil}^1)$ d'un isocristal de Dieudonné \add{admissible} sur $A_0 = A/pA$ (ou encore, sur $A_{0\,\mathrm{red}}$ \struck{\ill{}}) et d'une filtration en deux crans par modules localement libres de $(M_A)_p$, ayant le rang qu'il faut (et soumis éventuellement à d'autres conditions ouvertes \add{au sens Zar}, assurant que le couple envisagé provient d'un isogroupe de BT sur $A$), vers les faisceaux p-adiques constants tordus sur $U = \mathrm{Spec}(A_p)$, qui exprime le foncteur du théorème (VII~10.8.2) précédent.

\struck{4) Question des relations entre cohomologie cristalline et p-adique.}

\struck{Conjecture~7. On peut, pour tout anneau de valuation discrète complet $V$ d'inégales caractéristiques, de corps résiduel parfait (pour simplifier), de corps des fractions $L$, trouver un foncteur canoniques sur une sous-catégorie pleine convenable des couples $(M, \mathrm{Fil}^1)$ dans le Pb~5 vers les}
\note{ce dernier passage est entouré et biffé à la main de traits obliques ; il s'interrompt sur « vers les » ; à l'intérieur, la machine avait déjà biffé « Problème » (remplacé par « Conjecture »), « des fractions » et « envisagés »}

\page{52}
\note{en tête de page, au crayon, au-dessus d'un trait horizontal}
\[
  \left\{
  \begin{aligned}
    &\textstyle\sum_1^i \lambda'_j \leq \sum_1^i \lambda_j \\
    &\textstyle\sum_1^n \lambda'_j = \sum \lambda_j
  \end{aligned}
  \right.
  \qquad \textstyle\sum i h'_i
\]
\marginal{\uncertain{mieux de} partir des Hodge}

\emph{Sur un théorème de Tate.} (Devrait venir avec foncteur mystérieux\,!)

C'est le théorème suivant (Driebergen~….)\,:

\emph{Th.}\ (Tate)\,: Soit $S$ un schéma noethérien normal intègre dont le corps des fonctions \add{$k(\eta)$} est de car.\ zero, Alors le foncteur
\[
  \mathrm{BT}(S) \longrightarrow \mathrm{BT}(\eta)
\]
est pleinement fidèle.

Ainsi, un groupe de BT sur $S$ est connu \add{à isom unique près} quand on connait la représentation p-adique qu'il définit de $\pi_1(\eta) = \mathrm{Gal}(\overline{k(\eta)}/k(\eta))$. Prenant d'abord le cas où $S$ est le spectre d'un anneau de valuation discrète complet $V$ d'inégales car., de corps \struck{des} résiduel $k$ parfait, la question se pose donc

a) De caractériser les modules galoisiens qu'on peut obtenir (Tate donne des conditions nécessaires\,: les modules en question doivent être du type de «~Hodge-Tate~»)\,;

b) De trouver un foncteur «~explicite~» sur la catégorie de ces modules galoisiens, vers les F-cristaux sur $k$ $M$ muni d'une filtration de $M \otimes_W L$.
\note{« $\eta$ » est ajouté à la main partout dans ce paragraphe, dans les blancs du tapuscrit}

\add{\emph{Problème}} Le théorème de Tate suggère de plus que le «~foncteur mystérieux~» qui va des F-cristaux sur $k$ filtrés sur $L$, vers les $\underline{Q}_p$-modules galoisiens sur $L$, est pleinement fidèle, et qu'on doit pouvoir donc \struck{définir un foncteur} expliciter un foncteur quasi-inverse (non moins mystérieux), allant d'une sous-catégorie pleine convenable des $\underline{Q}_p$-modules galoisiens sur $L$ (sans doute contenue dans celle des modules galoisiens du type de Hodge-Tate) vers les F-cristaux sur $k$ filtrés sur $L$.

D'autre part, on conjecture que dans le th.\ de Tate ci-dessus, la restriction au cas où \struck{$S$} \add{$k(\eta) = L$} est de car.\ nulle est inutile, Le seul cas qui offre

\page{53}
un Pb.\ est celui où $k(\eta) = L$ est de car.\,p (cas d'égales caractéristiques). On peut alors penser à attaquer le Pb via la théorie de Dieudonné (supposant résolu le Pb~1) et via le

\emph{Problème} n\textsuperscript{o}~… Soit $S$ un schéma noethérien normal \struck{de car} intègre de car.\,p, de point générique $\eta$. Est-il vrai que le foncteur de restriction
\[
  \mathrm{Cris\ loc\ lib}(S) \longrightarrow \mathrm{Cris\ loc\ lib}(\eta)
\]
(ou du moins le foncteur analogue pour F-cristaux) est pleinement fidèle\,?

On est ramené facilement au cas où $S$ est le spectre d'un anneau de la forme $k[[t]]$, avec $k$ alg.\ clos, auquel cas l'énoncé s'explicite en termes de Modules avec connexions formelles~…. Il semble que si la réponse au Pb précédent est affirmative, \struck{ainsi que} et si on sait faire la théorie de Dieudonné sur un anneau \emph{parfait} (avec \emph{équivalence} de catégories alors on aura démonstration du th.\ de Tate en égales car.

\section{$\chi$ et $\chi_!$}
\note{titre écrit seul, de sa main, en tête d'une feuille autrement blanche (p.~54) ; les pages manuscrites qui suivent sont à l'encre bleue}

\page{55}
$k$ corps

$\mathcal{E}_{\chi_!}$ espaces topologiques \add{loc.\ compacts,} tels que $H^*_!(X,k)$ de dim finie sur $k$ \marginal{cohom.\ \struck{faisceaux} singulière}

$\mathcal{E}_{\chi}$ \quad --- \quad $H^*(X,k)$ \quad --- \marginal{cohom.\ singulière}

(a) \struck{\ill{}} si $X$ compact, $X \in \mathcal{E}_{\chi_!} \Longleftrightarrow X \in \mathcal{E}_{\chi}$, et lors
\[
  \chi_!(X) = \chi(X) .
\]
\add{trivial}

(b) si $Y$ fermé dans $X$ \add{loc.\ cpt}, alors si deux des espaces $X, Y, X-Y$ sont dans $\mathcal{E}_{\chi_!}$, la troisième aussi, et~…
\[
  \chi_!(X) = \chi_!(Y) + \chi_!(X-Y)
\]
(résulte de suite exacte de cohomologie à s-\uncertain{supports} \uncertain{cpts})

(c) si $X'$, $X''$, \struck{\ill{}} fermés dans $X$ de réunion $X$, alors si trois parmi les quatre espaces $X, X', X'', X' \cap X''$ sont dans $\mathcal{E}_{\chi}$, le quatrième aussi, et on
\[
  \chi(X) = \chi(X') + \chi(X'') - \chi(X' \cap X'')
\]
(résulte de suite exacte de Mayer-Vietoris)

\struck{(d)} [Mieux $X, X' \cap X'' \in \mathcal{E}_{\chi} \Longrightarrow X', X'' \in \mathcal{E}_{\chi}$]
\note{l'étiquette (d) de cette ligne est surchargée et biffée ; la ligne est entre crochets}

(d) si $X$ \struck{\ill{}} \add{«~dominé~»} homotopiquement par $X'$ ($X \overset{i}{\rightleftarrows} X'$, $p \circ i = \mathrm{id}_X$) alors $X' \in \mathcal{E}_{\chi} \Longrightarrow X \in \mathcal{E}_{\chi}$, et~…, $(X \sim X') \Longrightarrow (X \in \mathcal{E}_{\chi} \Longleftrightarrow X' \in \mathcal{E}_{\chi})$, et de plus \uncertain{on a}
\[
  \chi(X) = \chi(X')
\]

(e) \struck{\ill{} $X$, si $X$ connexe\,; un pt,} $\operatorname{Card} X \leq 1 \Longrightarrow X \in \mathcal{E}_{\chi}$, de plus on
\[
  \chi(\mathrm{pt}) = 1
\]
(de plus $\chi(\emptyset) = 0$, qui \uncertain{suit} c) p.ex) \quad (trivial)

\page{57}
\emph{Conclusions premières} (\uncertain{où} pour $\overline{\mathbb{R}} = [-\infty, +\infty]$)

\add{① Soient $X', X'' \in \mathcal{E}_{\chi}$ (resp.\ $\in \mathcal{E}_{\chi_!}$) $\Longleftrightarrow$ $X' \amalg X'' \in \mathcal{E}_{\chi}$ (resp.\ $\in \mathcal{E}_{\chi_!}$), \struck{résulte de c et e, et dans} $\chi(X) = \chi(X') + \chi(X'')$ (resp.~…) résulte de c) et e) ($\emptyset \in \mathcal{E}_{\chi}$) resp.\ de b)}
\note{cet énoncé ① est ajouté en haut à droite et encadré ; la numérotation de la page commence pourtant à ②}

② $\forall n \in \mathbb{N}$, \struck{\ill{}} $\mathbb{R}^n, \overline{\mathbb{R}}^n, B^n, S^n \in \mathcal{E}_{\chi} \cap \mathcal{E}_{\chi_!}$, de plus
\[
  \left\{
  \begin{aligned}
    &\chi(\mathbb{R}^n) = \chi(\overline{\mathbb{R}}^n) = \chi(B^n) = 1 \\
    &\phantom{\chi(\mathbb{R}^n) = {}} \chi_!(\overline{\mathbb{R}}^n) = \chi_!(B^n) \\
    &\chi_!(\mathbb{R}^n) = (-1)^n
  \end{aligned}
  \right.
  \qquad
  \chi(S^n) = \chi_!(S^n) = 1 + (-1)^n
\]
\note{des signes « $\parallel$ » placés au-dessus de $\chi_!(\overline{\mathbb{R}}^n)$ et $\chi_!(B^n)$ les égalent aux valeurs de la ligne précédente}

\emph{Dém.} $\mathbb{R}^n, \overline{\mathbb{R}}^n, B^n \sim \mathrm{pt}$, donc par d) et e) ils sont $\in \mathcal{E}_{\chi}$\,; comme $\overline{\mathbb{R}}^n$, $B^n$ compacts, par b) ils sont aussi $\in \mathcal{E}_{\chi_!}$. On a $S^n = S^n_+ \cup S^n_-$ et $S^n_+ \cap S^n_- \simeq S^{n-1}$, \add{or $S^n_+ \simeq S^n_- \simeq B^n$ donc $\in \mathcal{E}_{\chi}$, donc il \uncertain{conclut}} de c) par réc.\ sur $n$ par $S^n \in \mathcal{E}_{\chi}$ (\emph{NB} $S^0 = \{\mathrm{pt}\} \amalg \{\mathrm{pt}\}$, et $\chi(S^0) = 1 + 1 = 2$), et on trouve
\[
  \chi(S^n) = \chi(S^n_+) + \chi(S^n_-) - \chi(S^{n-1}) = 2 - \chi(S^{n-1}) ,
\]
d'où par récurrence sur $n$
\[
  \chi(S^n) = 1 + (-1)^n .
\]
Comme $S^n$ cpt, $S^n \in \mathcal{E}_{\chi_!}$, et $\chi_!(S^n) = \chi(S^n)$ par (a). Enfin, $\mathbb{R}^n \simeq S^n - \mathrm{pt}$, donc par (b) $\mathbb{R}^n \in \mathcal{E}_{\chi_!}$ et
\[
  \chi(\mathbb{R}^n) = \chi(S^n) - \chi(\mathrm{pt}) = (1 + (-1)^n) - 1 = (-1)^n ,
\]
\note{il écrit ici $\chi$ et non $\chi_!$, bien que ce soit $\chi_!(\mathbb{R}^n)$ que (b) calcule}
cqfd.

③ Soit $X_{-1} = \emptyset \subset X_0 \subset X_1 \subset \cdots \subset X_n = X$ filtration de l'espace $X$, tel que $\forall\, 0 \leq i \leq n$, $X_i - X_{i-1}$ \struck{\ill{}} a un \uncertain{ens} fini $f_i$ de comp.\ connexes \uncertain{toutes} homéom.\ à $\mathbb{R}^i$ (i.e.\ $X_i - X_{i-1} \simeq f_i \times \mathbb{R}^i$). Alors $X \in \mathcal{E}_{\chi_!}$, et~…
\[
  \chi_!(X) = \sum (-1)^i f_i
\]
(donc si $X$ est cpt~…, $X \in \mathcal{E}_{\chi}$ et $\chi(X) = \sum (-1)^i f_i$~…

\page{59}
\note{les points (f) et (g) continuent la liste (a)–(e) de la p.~55}

(f) $X, Y \in \mathcal{E}_{\chi_!} \Longrightarrow X \times Y \in \mathcal{E}_{\chi_!}$ et~… \quad $\chi_!(X \times Y) = \chi_!(X)\chi_!(Y)$

$X, Y \in \mathcal{E}_{\chi} \Longrightarrow X \times Y \in \mathcal{E}_{\chi}$ et~… \quad $\chi(X \times Y) = \chi(X)\chi(Y)$.

(La première assertion résulte de Künneth pour $H^*_!(X \times Y, k)$, la deuxième de Künneth pour $H_*(X \times Y, k)$, dont le dual est $H^*(X \times Y, k)$, ce qui permet de tenir l'\uncertain{hypothèse} \uncertain{dans} $\mathcal{E}_{\chi}$ et la valeur de $\chi$~…)

(g) Soit $f : X \to Y$ fibration loc.\ triviale de fibre $F$. Si $Y, F \in \mathcal{E}_{\chi_!}$\,; alors $X \in \mathcal{E}_{\chi_!}$ et~…
\[
  \chi_!(X) = \chi_!(Y)\chi_!(F) .
\]
Énoncé analogue pour \add{$Y, F \in \mathcal{E}_{\chi}$} \struck{$\mathcal{E}_{\chi}$} et \struck{$\chi(X)$ \ill{}}
\[
  \chi(X) = \chi(Y)\chi(F)
\]
\marginal{avec \uncertain{lign\ldots} \uncertain{générales} \ill{}}

\emph{Dém.} On a
\[
  H^*_!(X) \Longleftarrow E_2^{pq} = H^p_!(Y, \underbrace{R^q f_!(k_X)}_{\substack{\text{syst.\ local sur } Y \\ \text{à fibres} \sim H^q_!(F)}})
\]
donc \emph{sans réserve} puisque on sait que $E_2^{**}$ de dim finie, on trouve $X \in \mathcal{E}_{\chi}$ et
\[
\begin{aligned}
  \chi(X) &= \sum (-1)^{p+q} \operatorname{rg}_k H^p_!(Y, R^q f_!(k_X)) \\
  &= \sum_q (-1)^q \underbrace{\sum_p (-1)^p \operatorname{rg}_k H^p_!(Y, R^q f_!(k_X))}_{b^q_!(F)\chi_!(Y)} \\
  &= \chi_!(Y) \sum_q (-1)^q b^q_!(F) = \chi_!(Y)\chi_!(F)
\end{aligned}
\]
\note{il écrit $X \in \mathcal{E}_{\chi}$ et $\chi(X)$ là où l'énoncé porte sur $\mathcal{E}_{\chi_!}$ et $\chi_!(X)$ ; sous la somme, « $\sum_q$ » est surmonté d'une marque « $p,q$ »}

Il faut prouver

\emph{Lemme}\,: Soit $\mathcal{R}$ syst.\ local de $k$-vectoriels, sur \struck{$X$ loc.\ cpt,} $Y \in \mathcal{E}_{\chi_!}$, à fibres de rang \add{fini} $b$. Alors
\[
  \chi_!(Y, \mathcal{R}) \overset{\mathrm{déf}}{=} \sum_p (-1)^p H^p(Y, \mathcal{R})
\]
existe (i.e.\ $\sum_p H^p(Y, \mathcal{R})$ de dim finie sur $k$) et~…
\[
  \chi_!(Y, \mathcal{R}) = b\,\chi_!(Y)
\]

Marche tout au moins si $Y$ \struck{triangulable} \add{(gr) admet filtration finie} avec $Y_i - Y_{i-1}$ sommes finies de cellules \add{$Z_{ij}$} $\simeq \mathbb{R}^{d_{ij}}$, car on aura
\[
  \chi_!(Y, \mathcal{R}) = \sum_{ij} \chi_!(Z_{ij}, \mathcal{R}) \underset{(\text{cellules } Z_{ij})}{=} \sum_{ij} b\,\chi_!(Z_{ij}) = b \sum \chi_!(Z_{ij}) = b\,\chi_!(Y) .
\]
\note{« (gr) » est cerclé ; « sommes finies de cellules » est une lecture de la ligne}

\end{document}
